% Robustness and validity: each paragraph tests one way the results could be wrong.
\section{Robustness and Validity}
\label{sec:robustness}

Each check below tests one way in which the results could be wrong. All of them are generated by the released code;
the full tables are in Appendix~\ref{app:tables}.

\paragraph{Did each model receive the intended prompt?}
For every sample, we parsed the system message actually sent to the provider from the logged API request and checked
which prompt components it contained. Of \RNSamples{} samples, \RNWrongPrompt{} had a wrong system prompt and
\RNJudgeLeak{} had a judge call carrying any condition prompt.

\paragraph{Multiplicity.}
Correcting across all twelve primary tests instead of within each contrast leaves both main effects significant: the
HumaneBench drop at $p=\RHbCevClassicHolmAll$ and the AI-risk gain at $p=\RAirClassiccevLenmatchHolmAll$. The
HumanAgencyBench drop does not reach significance ($p=\RHabClassiccevLenmatchHolmAll$).

\paragraph{Does one model drive an effect?}
Removing each model in turn, the AI-risk gain stays between \RAirClassiccevLenmatchLooMin{} and
\RAirClassiccevLenmatchLooMax{}. The HumaneBench drop stays negative (\RHbCevClassicLooMin{} to
\RHbCevClassicLooMax{}) but depends mostly on Llama. Without Llama it is \RHbCevClassicWithoutLlamaDiff{} (Holm
$p=\RHbCevClassicWithoutLlamaHolm$). See Appendix Table~\ref{tab:loo}.

\paragraph{Missing data.}
\RHbUnscorableShare{} of HumaneBench samples are unscorable, because its rubric lets the judge mark the targeted
principle as not applicable. The pattern is not random: items unscorable under \cond{cev} had lower scores under
\cond{classic}. We therefore re-estimated the HumaneBench contrasts two ways:
\begin{itemize}
  \item on complete cases (items scored in all five conditions): \RHbMissCompleteCevClassicDiff{},
  $p=\RHbMissCompleteCevClassicP$;
  \item imputing unscorable samples at three fixed values: \RHbMissImputeCevClassicMin{} to
  \RHbMissImputeCevClassicMax{}, all $p\le\RHbMissImputeCevClassicPMax$.
\end{itemize}
The conclusion is unchanged (Appendix Table~\ref{tab:missing}).

\paragraph{Provider routing.}
OpenRouter served the open-weight models from several backends (\RQwenProviders{} for Qwen, \RLlamaProviders{} for
Llama), and because the conditions ran one after another, the mix differed between conditions. After removing each
evaluation$\times$condition$\times$model mean, no backend with at least 50 samples deviates by more than
\RProviderResidMax{} in aligned score (Appendix Table~\ref{tab:providers}). This cannot account for the effects
above. Providers were not pinned, however.

\paragraph{Where do the losses come from?}
Under the two CEV prompts, \RMentionShareMin--\RMentionShareMax{} of answers explicitly invoke CEV or ``extrapolated
volition'', reaching \RMentionShareCevOnlyMax{} under \cond{cev} alone. These answers score much lower than the rest:
\begin{itemize}
  \item HumanAgencyBench: \RHabMentionAligned{} against \RHabNoMentionAligned{};
  \item HumaneBench: \RHbMentionAligned{} against \RHbNoMentionAligned{};
  \item Agentic Misalignment: \RAmMentionAligned{} against \RAmNoMentionAligned{} (Table~\ref{tab:mentions}).
\end{itemize}
Whether an answer mentions CEV is itself caused by the prompt, so this split is descriptive. It fits two
explanations: models that reason ``as CEV'' behave differently, or the judge penalises the mention.

\begin{table}[t]
\centering\small
% Generated by cev_alignment.paper_assets. Do not edit.
\begin{tabular}{lcccc}
\toprule
Eval & Share invoking CEV & Aligned (invoking) & Aligned (not invoking) & $n$ invoking \\
\midrule
Agentic Misalignment & 15\% & 0.56 & 0.97 & 18 \\
AI-risk (MWE) & 16\% & 0.78 & 0.79 & 118 \\
DarkBench & 15\% & 0.81 & 0.80 & 73 \\
HumanAgencyBench & 12\% & 0.18 & 0.46 & 57 \\
HumaneBench & 16\% & 0.72 & 0.87 & 62 \\
InstrumentalEval & 1\% & 1.00 & 0.84 & 3 \\
\bottomrule
\end{tabular}

\caption{Answers under \cond{cev} and \cond{classic+cev} that explicitly invoke CEV, and the mean aligned score of
answers that do and do not invoke it. Descriptive: invoking CEV is an outcome of the prompt.}
\label{tab:mentions}
\end{table}

\paragraph{Judge-blinding test.}
To separate the two explanations, we re-judged all \RBlindN{} CEV-invoking answers twice, with the same scorer and
judge: once as written, and once with every sentence that mentions CEV deleted. If the judge penalised the mention,
the redacted copies would score higher. They do not (Table~\ref{tab:blinding}). HumanAgencyBench changes by
\RBlindHabDiff{} ($p=\RBlindHabP$), HumaneBench by \RBlindHbDiff{} ($p=\RBlindHbP$), and DarkBench by \RBlindDbDiff{}
($p=\RBlindDbP$). Redaction can also delete substantive content, so these differences are an upper bound on any
judge reaction. The judge is deterministic: re-judging gave identical scores for \RRetestAgree{} of \RRetestN{}
answers. The lower scores therefore reflect what CEV-invoking answers do.

\begin{table}[t]
\centering\small
% Generated by cev_alignment.paper_assets. Do not edit.
\begin{tabular}{lccccc}
\toprule
Eval & $n$ & Original & Redacted & $\Delta$ [95\% CI] & $p$ \\
\midrule
Agentic Misalignment & 18 & 0.556 & 0.556 & \ensuremath{+0.000} [\ensuremath{+0.000}, \ensuremath{+0.000}] & 1 \\
DarkBench & 71 & 0.803 & 0.789 & \ensuremath{-0.014} [\ensuremath{-0.099}, \ensuremath{+0.070}] & 1 \\
HumanAgencyBench & 54 & 0.170 & 0.150 & \ensuremath{-0.020} [\ensuremath{-0.054}, \ensuremath{+0.011}] & 0.23 \\
HumaneBench & 43 & 0.733 & 0.733 & \ensuremath{+0.000} [\ensuremath{-0.052}, \ensuremath{+0.052}] & 0.87 \\
InstrumentalEval & 3 & 1.000 & 1.000 & \ensuremath{+0.000} [\ensuremath{+0.000}, \ensuremath{+0.000}] & 1 \\
\bottomrule
\end{tabular}

\caption{Judge-blinding test. Mean aligned score of CEV-invoking answers as written (Original) and with the
CEV-mentioning sentences removed (Redacted); $\Delta$ = redacted $-$ original. A positive $\Delta$ would mean the
judge penalises the mention itself.}
\label{tab:blinding}
\end{table}

\paragraph{Is the judge valid?}
We re-judged \RValN{} stored HumanAgencyBench responses from \RValModels{} different subject models with our judge,
using the benchmark's rubric, and compared the results with the benchmark's original judge scores. Agreement is
strong:
\begin{itemize}
  \item Spearman $\rho=\RValRho$ and Pearson $r=\RValPearson$;
  \item mean absolute error \RValMae{} on the 0--10 scale;
  \item Cohen's $\kappa=\RValKappa$ on whether a response earns a perfect score.
\end{itemize}
Our judge is slightly stricter (mean difference \RValBias{}). Appendix Table~\ref{tab:validity} breaks this down by
dimension. We did not validate the judge on HumaneBench or DarkBench.
